% ==============================================================================
%  3.7  Medidas de la relación lineal entre dos variables.
%       Coeficiente de determinación y correlación lineal
% ==============================================================================
\section{Medidas de la relación lineal: determinación y correlación}
\label{sec:t3-correlacion}

\subsection{El coeficiente de determinación}

En la descomposición de la variabilidad (\cref{teo:t3-sst}), $SST$ mide la incertidumbre en la
predicción de $Y$ sin usar la información de $X$, y $SSE$ la incertidumbre al usarla. La
reducción de la variabilidad, $SST-SSE=SSR$, en proporción a la variabilidad total, mide cuánto
se reduce la incertidumbre al usar $X$ para predecir $Y$.

\begin{definicion}[Coeficiente de determinación][def:R2]
  \[
    R^2 = \frac{SSR}{SST} = 1-\frac{SSE}{SST}.
  \]
\end{definicion}

Como $0\le SSE\le SST$, se tiene $0\le R^2\le1$:
\begin{itemize}
  \item $R^2=0$ cuando $SSR=0$: la recta ajustada es horizontal ($\est{\beta}_1=0$,
    $\est{Y}=\media{Y}$) y no hay relación lineal entre $X$ e $Y$;
  \item $R^2=1$ cuando $SSE=0$: todas las observaciones coinciden con sus valores ajustados, y
    $X$ predice perfectamente a $Y$.
\end{itemize}
En la práctica se obtienen valores intermedios. Cuanto mayor es $R^2$, mayor es la reducción de
la variabilidad total de $Y$ al usar $X$ para predecirla.

\paragraph{Propiedades de $R^2$.}
\begin{itemize}
  \item $0\le R^2\le1$ y es adimensional.
  \item Es el cuadrado del coeficiente de correlación de Pearson de la muestra
    $\{(X_i,Y_i)\}_{i=1}^n$, y también del de los pares $\{(\est{Y}_i,Y_i)\}_{i=1}^n$.
  \item Cuanto mayor es $R^2$, mayor es la capacidad del regresor $X$ para predecir la respuesta
    $Y$, menor es $SSE$ y más cerca están los puntos de la recta ajustada.
  \item Toma el mismo valor si se usa $X$ para predecir $Y$ o $Y$ para predecir $X$.
\end{itemize}
En el \cref{ej:t3-anova-arboles}, $R^2=3085.789/3744.358=0.8241$: el diámetro del tronco explica el
82.4\,\% de la variabilidad del volumen de madera.

\subsection{El coeficiente de correlación lineal}

Cuando las dos variables son aleatorias, la medida habitual de la relación lineal entre ellas es
el coeficiente de correlación de Pearson muestral (\cref{def:pearson}), que en función de las
sumas de cuadrados y de la pendiente estimada se escribe
\[
  r_{XY} = \frac{\sumi (x_i-\media{X})(y_i-\media{Y})}
                {\sqrt{\sumi (x_i-\media{X})^2\sumi (y_i-\media{Y})^2}}
  = \frac{SS_{XY}}{\sqrt{\SSX\,SS_Y}} = \est{\beta}_1\sqrt{\frac{\SSX}{SS_Y}} .
\]

\begin{proposicion}[Relación entre $R^2$ y $r_{XY}$][prop:R2-r2]
  En la regresión lineal simple, $R^2=r_{XY}^2$.
\end{proposicion}

\begin{demostracion}
  Usando $SSR=\est{\beta}_1^2\SSX$ (\cref{sec:t3-anova}) y $SS_Y=SST$,
  \[
    r_{XY}^2 = \est{\beta}_1^2\,\frac{\SSX}{SS_Y} = \frac{SSR}{SST} = 1-\frac{SSE}{SST} = R^2 .
  \]
\end{demostracion}

\paragraph{Propiedades de $r_{XY}$.}
\begin{itemize}
  \item $-1\le r_{XY}\le1$.
  \item $|r_{XY}|$ mide la fuerza de la asociación lineal: cuanto mayor es, más fuerte es la
    asociación lineal entre $X$ e $Y$.
  \item $r_{XY}=0$ indica que no hay asociación lineal; $r_{XY}=1$ ($-1$), asociación lineal
    perfecta positiva (negativa): los puntos están sobre una recta de pendiente positiva
    (negativa).
  \item El signo indica asociación positiva ($r_{XY}>0$) o negativa ($r_{XY}<0$).
  \item No depende de las unidades de medida de las variables.
  \item Los papeles de $X$ e $Y$ son simétricos en su cálculo.
  \item Es muy sensible a observaciones atípicas.
\end{itemize}

\begin{nota}[Regresión y correlación]
  Son conceptos relacionados que no deben confundirse. El modelo de regresión describe la
  relación lineal entre una variable independiente $X$, fija, y una respuesta $Y$, aleatoria,
  cuyo comportamiento se quiere explicar o estimar. La correlación describe la fuerza de la
  relación lineal entre dos variables, ambas aleatorias.
\end{nota}

\subsection{Inferencias sobre el coeficiente de correlación poblacional}

Las inferencias sobre el coeficiente de correlación poblacional $\rho_{XY}$ se basan en la
distribución en el muestreo de $r_{XY}$.
\begin{itemize}
  \item Para $H_0\colon\rho_{XY}=0$ frente a $H_1\colon\rho_{XY}\neq0$ se utiliza que
    \[
      r_{XY}\,\frac{\sqrt{n-2}}{\sqrt{1-r_{XY}^2}} \bajoHo \tdist{n-2} .
    \]
  \item Para $\rho_{XY}\neq0$ la distribución de $r_{XY}$ es compleja y se usan aproximaciones.
    La más habitual es la \textbf{transformación $Z$ de Fisher}: para $n$ grande,
    \[
      z' = \frac12\ln\left(\frac{1+r_{XY}}{1-r_{XY}}\right)
      \ \approx\ \Normal\!\left(\frac12\ln\left(\frac{1+\rho_{XY}}{1-\rho_{XY}}\right),\ \frac{1}{n-3}\right).
    \]
\end{itemize}
